\section{Le réalisme d'une monarchie chrétienne}

\subsection*{Introduction}

\subsection{Questionnements sur le modèle monarchique}

\subsubsection{Acceptation et mise en place d'un souverain unique}

L'analyse de la pensée politique d'Ambroise de Milan montre qu'il ne se positionne pas comme un révolutionnaire institutionnel, mais plutôt comme un penseur capable d'encadrer un gouvernement et d'apporter des fondamentaux moraux afin que ce dernier agisse pour le bien commun. Dans cet objectif, il ne délaisse donc pas le système monarchique, qu'il voit comme un système politique envisageable pour l'idéal chrétien s'il est correctement guidé. Bien qu'élogieux à l'égard des idéaux républicains, il est conscient de la réalité contemporaine du système impérial et de la nécessité d'un empereur unique à la tête de la société romaine. C'est ce que souligne Francis Dvornik sur la capacité d'Ambroise à accepter la monarchie dans sa théorie politique : «~Il [Ambroise] considère que la perfection républicaine peut être atteinte sous une sage monarchie, à condition que le monarque et les sujets comprennent leurs devoirs.\autocite[676]{dvornik_political}~» Pour argumenter cette idée, Ambroise s'appuie de nouveau sur son observation de l'ordre naturel. Il ne s'arrête pas juste aux sociétés animales fonctionnant de manière collégiale, mais prend aussi exemple sur les monarchies naturelles œuvrant pour le bien commun afin d'en tirer un modèle applicable à l'Empire romain, qui peut s'accorder avec son idéal chrétien.

\bigskip

C'est pour cette raison qu'il parle longuement des abeilles dans l'\latin{Exameron}, qui vivent sous la domination d'un roi\footnote{La science antique n'avait pas encore pris connaissance du fait que la ruche est dominée par une reine et non par un roi, ce qui explique d'autant plus le choix d'Ambroise de faire des abeilles un exemple pour sa théorie politique.} :

\bigskip

\begin{quote}
    «~Elles s'instituent elles-mêmes un roi, elles créent elles-mêmes leurs peuples et, bien qu’elles soient placées sous l'autorité d'un roi, elles demeurent pourtant libres. Car elles détiennent à la fois le privilège du jugement et le sentiment d'une dévotion fidèle ; puisqu'elles aiment celui qu'elles ont elles-mêmes substitué, et l'honorent avec une attention si scrupuleuse par leur grand nombre.\footnote{« \latin{Ipsae sibi regem ordinant, ipsae populos creant et licet positae sub rege sunt tamen liberae. nam et praerogatiuam iudicii tenent et fidae deuotionis affectum, quia et tanquam a se substitutum diligunt et tanto honorant examine}.~». \cite[\nopp V, 21, 68]{ambroise_csel_32_1}. Traduction personnelle.}~»
\end{quote}

\bigskip

À travers cet exemple naturel, Ambroise donne un sens et une légitimité à la royauté. Il démontre que la présence d'un souverain unique n'est pas incompatible avec la liberté et l'intérêt général. Mieux encore, l'autorité monarchique, si elle répond aux besoins de la communauté, permet un respect mutuel et une consolidation de la société. L'idéal du gouvernement chrétien se retrouve dans les principes de liberté et de dévotion envers autrui. On comprend bien que la plus grande crainte de l'évêque reste la présence d'un pouvoir coercitif et usant de la violence envers ses citoyens, ce que la monarchie favorise par essence. Son acceptation de celle-ci vient donc du respect mutuel avec les citoyens et de l'encadrement moral, ce que nous développerons dans la suite de ce chapitre.

\bigskip

Mais l'analyse d'Ambroise sur la monarchie des abeilles prend aussi sa justification par le fait qu'elles représentent pour lui le paroxysme de la société de partage et d'intérêt commun, alors même qu'elles obéissent à un roi. Dans son paragraphe introduisant cette société, il appuie sur la notion de mise en commun :

\bigskip

\begin{quote}
    «~Ce sont là de grandes choses, mais combien elles sont supérieures chez les abeilles ! Seules parmi toutes les espèces vivantes, elles ont une progéniture commune à toutes, elles habitent toutes une seule demeure, et sont enfermées dans le seuil d'une seule patrie. En commun pour toutes est le travail, commune est la nourriture, commune à toutes est l'activité, communs sont l'usage et le fruit [...].\footnote{« \latin{Magna haec, sed quanto in apibus praestantiora, quae solae in omni genere animantium communem omnibus subolem habent, unam omnes incolunt mansionem, unius patriae clauduntur limine. in commune omnibus labor, communis cibus, communis omnibus operatio, communis usus et fructus est [...]}.~». \cite[\nopp V, 21, 67]{ambroise_csel_32_1}. Traduction personnelle.}~»
\end{quote}

\bigskip

En martelant le terme \latin{communis}, Ambroise insiste sur ce qui est le plus important de sa théorie politique : le partage et l'intérêt général. Ce qui compte dans la théorie ambrosienne, ce n'est pas le système politique, mais l'engagement des citoyens et des dirigeants à agir pour la réussite de la communauté.

\bigskip

Mais l'acceptation d'un souverain unique par Ambroise ne se fait pas seulement par la réussite du monde animal : elle s'appuie également sur une théologie de l'Histoire romaine faisant un lien entre la diffusion du christianisme et l'arrivée au pouvoir d'Auguste. Il s'associe ainsi à une pensée traditionnelle qui légitime la présence d'un monarque comme un outil ayant facilité le développement des communautés chrétiennes. Cette justification du système impérial apparaît dans son commentaire du Psaume 45 :

\bigskip

\begin{quote}
    «~C’est pourquoi, lassé des guerres civiles, l’empire romain fut confié à Jules Auguste, et les conflits internes furent apaisés. Cela permit également l’envoi légitime des Apôtres dans le monde, le Seigneur Jésus disant : "Allez, de toutes les nations faites des disciples."\footnote{« \latin{Unde factum est, ut taedio bellorum civilium, Julio Augusto Romanum deferretur imperium : et ita praelia intestina sedata sunt. Hoc autem eo profecit, ut recte per totum orbem apostoli mitterantur dicente Domino Jesu : Euntes docete omnes gentes}.~». \cite[\nopp 45, 21]{ambroise_csel_64}.}~»
\end{quote}

\bigskip

Il y a chez Ambroise cette idée que le système impérial peut se légitimer dans une idéologie chrétienne par le simple fait que c'est par lui que la foi chrétienne a pu atteindre son plein essor. Le fait même d'avoir un roi à la tête de la société n'est donc pas forcément vu comme quelque chose de négatif entravant forcément un idéal politique. Évidemment, on note une justification de ses propos par la présence d'un passage du Nouveau Testament qui fait du christianisme une religion à vocation universaliste, ce qui est largement facilité par la présence d'une seule autorité tout autour de la Méditerranée. Cette pensée est d'ailleurs à replacer dans une tradition apologétique plus large que juste l'évêque de Milan. Il rejoint de très près, y compris dans les termes utilisés, les thèses d'Origène un siècle et demi plus tôt. Dans son \latin{Contre Celse}, ce dernier établissait déjà ce synchronisme entre Jésus et Auguste, entre l'Église et l'Empire :

\begin{quote}
    «~Dieu préparait les nations à recevoir son enseignement, en les faisant toutes obéir au seul empereur de Rome, et en empêchant que l’isolement des nations dû à la pluralité des royautés ne rendît plus difficile aux apôtres l’exécution de l’ordre du Christ. [...] Il est manifeste que Jésus est né sous le règne d’Auguste qui avait pour ainsi dire ramené à une masse uniforme, grâce à sa royauté unique, la plupart des royaumes de la terre.\footnote{« \gr{εὐτρεπίζοντος τοῦ θεοῦ τῇ διδασκαλίᾳ αὐτοῦ τὰ ἔθνη, ἵν' ὑπὸ ἕνα γένηται τῶν Ῥωμαίων βασιλέα, καὶ μὴ διὰ τὸ προφάσει τῶν πολλῶν βασιλειῶν ἄμικτον τῶν ἐθνῶν πρὸς ἄλληλα χαλεπώτερον γένηται τοῖς ἀποστόλοις τοῦ Ἰησοῦ τὸ ποιῆσαι ὅπερ προσέταξεν αὐτοῖς ὁ Ἰησοῦς εἰπών. [...] Καὶ σαφές γε ὅτι κατὰ τὴν Αὐγούστου βασιλείαν ὁ Ἰησοῦς γεγέννηται, τοῦ, ἵν' οὕτως ὀνομάσω, ὁμαλίσαντος διὰ μιᾶς βασιλείας τοὺς πολλοὺς τῶν ἐπὶ γῆς.}.~». \cite[\nopp II, 30]{origene_celse}.}~»
\end{quote}

\bigskip

Cette conception de l'Histoire reliant l'autorité impériale aux nouvelles institutions de l'Église légitime la présence d'un souverain unique, en rappelant naturellement l'argument monarchique du «~Un Dieu, un Empire, un empereur~»\autocite[681]{dvornik_political}. L'unification géopolitique apportée par les victoires d'Auguste a agi comme un outil facilitant la tâche du Christ. Ambroise reste donc dans la continuité d'une pensée qui voit l'Empire comme un instrument structurel indispensable à l'Église, une vision qui sera remise en question avec Augustin après lui, lorsqu'il se questionnera sur la capacité de l'Église à survivre à un potentiel effondrement de l'Empire romain.

\bigskip

En définitive, si la théorie politique ambrosienne favorise la présence d'un système républicain, elle accepte pleinement la mise en place du modèle monarchique, qui semble inévitable dans la société romaine. Derrière sa légitimation par les exemples de la nature ou par la réussite de la diffusion du christianisme, cette monarchie est décrite comme un moindre mal politique et moral, si le dirigeant est bien encadré par la morale chrétienne.

\bigskip

\subsubsection{Un système monarchique qui nécessite un respect mutuel}

Si le système monarchique est accepté par Ambroise comme une réalité contemporaine difficile à dépasser, sa viabilité reste pour autant dépendante des actions et volontés du souverain. Dans son objectif de théoriser le gouvernement romain, il ne laisse pas le pouvoir de l'empereur comme une finalité en soi, mais comme un aspect pouvant constamment s'améliorer pour aboutir à l'idéal chrétien du bon gouvernement. La réussite de la monarchie ne peut reposer sur la seule force coercitive des institutions impériales. Pour que ce gouvernement s'accorde avec les réflexions ambrosiennes, il doit s'incarner autour d'un contrat moral qui parvient à lier de manière réciproque les citoyens et les détenteurs du pouvoir.

Pour trouver les fondements de la pensée d'Ambroise à ce sujet, il faut évidemment se plonger dans le \latin{De Officiis}, qui lui sert d'exutoire moral et de support pour poser toute sa réflexion politique et ses volontés d'encadrement des dirigeants. Il évoque ainsi à son clergé ce qu'il considère comme les causes de la chute ou de la réussite d'une société. Il écarte l'idée que les menaces envers la monarchie ne proviennent que de l'extérieur pour, au contraire, intimider les souverains en rappelant que l'effondrement de l'intérieur est plus commun que celui venant des défaites militaires :

\bigskip

\begin{quote}
    «~Nous lisons en effet, non seulement en ce qui concerne les particuliers, mais aussi les rois eux-mêmes, combien fut profitable l’aisance d’une séduisante affabilité, ou combien furent dommageables l'orgueil et la hauteur des paroles, au point d’ébranler les royaumes eux-mêmes et de briser la puissance. Donc si quelque roi par sa sagesse, sa façon d’agir, son administration, l’accomplissement de ses devoirs, gagne la faveur du peuple, ou si quelque roi se présente au danger pour l’ensemble de la population, ce n’est pas douteux : un tel amour refluera de la population vers lui, que le peuple fera passer le salut du roi avant son intérêt.\footnote{« \latin{texte en latin}.~». \cite[\nopp II, 30]{ambroise_devoirs_1}.}~»
\end{quote}

\bigskip

Pour Ambroise, des notions déjà rapidement évoquées sont au cœur de l'échec des monarchies et de l'impossibilité d'accéder à son idéal : l'orgueil et l'abus de puissance. Ce sont ces fautes qui amènent la corruption de la société et le développement d'un gouvernement coercitif n'agissant plus pour le bien commun. Cette citation révèle l'exigence de la théorie politique ambrosienne. Il apporte d'ailleurs des éléments de réflexion assez précis en dressant une liste des conditions de la réussite de la monarchie : la sagesse, l'administration et l'accomplissement des devoirs. Le pouvoir du monarque ne vient pas seulement de son titre et de ses droits institutionnels : il doit le mériter au quotidien par ses actions pour être certain d'en conserver la légitimité, et ainsi éviter de tomber dans l'obligation d'un maintien de l'ordre par la violence à cause de l'absence de respect du pouvoir par les citoyens. Puis Ambroise continue avec des termes forts pour s'adresser au souverain, qui doit «~se présenter au danger pour l'ensemble de la population~». Bien que la société romaine ne soit pas une république comme peut l'être celle des grues, le partage des devoirs et l'implication des détenteurs du pouvoir envers les citoyens n'en restent pas moins importants. Une monarchie qui s'impose comme un gouvernement juste et idéal, c'est une monarchie capable de se sacrifier au service de la société. Mais c'est finalement la fin de la citation qui résume le mieux la pensée de l'évêque : la réciprocité dans l'engagement envers l'autre. Si le roi gouverne avec douceur et justesse, alors le peuple le lui rendra par le respect et le suivi général.

\bigskip

Ce respect demandé au monarque envers son peuple ne reste pas simplement un concept philosophique vague. Comme souvent avec Ambroise, sa proximité avec les empereurs de son temps lui permet d'exploiter pleinement sa réflexion théorique et de la faire appliquer à travers les nombreuses correspondances qu'il entretient avec le cœur du pouvoir. On le retrouve par exemple dans sa lettre 75, adressée à l'empereur Valentinien II dans le cadre du conflit sur la possession des basiliques à Milan\footnote{Voir la sous-partie [...] du chapitre 3}. Il illustre de façon concrète comment la monarchie doit s'auto-encadrer :

\bigskip

\begin{quote}
    «~Et lorsque vous avez posé cette règle pour les autres, vous l'avez posée aussi pour vous-même ; car les lois que l'Empereur fait, il doit être le premier à les respecter.\footnote{« \latin{texte en latin}.~». \cite[\nopp 75, 9]{ambroise_csel_82_3}.}~»
\end{quote}

\bigskip

Cette consigne adressée avec fermeté est l'exemple parfait de ce que l'évêque entend par le respect mutuel entraînant l'encadrement de la monarchie à travers ses propres actions. Pour que le monarque se mette au service de sa population comme Ambroise l'indique dans le \latin{De Officiis}, il faut commencer par se limiter aux règles qu'il met lui-même en place. Si le souverain refuse cet aspect et se considère au-dessus des lois, non seulement il tombe dans la faute vis-à-vis de la foi chrétienne, mais en plus il entraîne avec lui l'ensemble de la société qui passe sous le joug de la domination par la force. L'acceptation du système impérial par Ambroise est donc bien loin d'une validation de la puissance absolue qui, comme elle n'est pas partagée, n'est pas non plus limitée. L'exemplarité juridique de l'empereur se présente comme le fondement du respect envers la communauté qu'il dirige et comme un point intransigeant de la vision politique d'Ambroise.

\bigskip

C'est précisément cet auto-encadrement du monarque, et surtout la mise en retrait de l'orgueil et de la domination au profit de la justice, qui permet la pleine acceptation d'un souverain unique par le peuple. Ambroise veut valoriser le bon gouvernement autrement que par le succès matériel de la société ; il montre alors dans son \latin{Exameron} que lorsque le pouvoir se montre juste, il obtient un retour réciproque de la population, qui se doit d'aller loin dans le service à l'Empire. Son allégorie politique avec les abeilles lui permet de mettre en avant cette réflexion :

\bigskip

\begin{quote}
    «~Mais même celles qui n'ont pas obéi aux lois du roi se punissent elles-mêmes par une condamnation repentante, au point de mourir de la blessure de leur propre aiguillon. C'est ce que, dit-on, les peuples des Perses observent encore aujourd'hui : pour le prix de leur faute, ils exécutent eux-mêmes sur leur propre personne la sentence de mort. Ainsi, aucun peuple n'observe son roi avec autant de révérence et de dévouement que les abeilles, [...] de sorte qu'aucune d'elles n'ose sortir de ses demeures ni s'avancer vers quelque pâturage, si le roi n'est sorti le premier et n'a revendiqué pour lui la direction du vol.\footnote{« \latin{Sed etiam illae quae non obtemperauerint legibus regis paenitenti condemnatione se multant, ut inmoriantur aculei sui uulneri. quod Persarum populi hodieque seruare dicuntur, ut pro commissi pretio ipsi in se propriae mortis exequantur sententiam. itaque nulli sic regem, [...] tanta quanta apes reuerentia deuotionis obseruant, ut nullae de domibus exire audeant, non in aliquos prodire pastus, nisi rex fuerit primo egressus et uolatus sibi uindicauerit principatum}.~». \cite[\nopp V, 21, 68]{ambroise_csel_32_1}. Traduction personnelle.}~»
\end{quote}

\bigskip

Cette nouvelle citation de son commentaire des six jours de la Création nous permet de voir une continuité politique cohérente chez Ambroise vis-à-vis des éléments que nous avons précédemment analysés : l'implication du peuple dans le maintien de la mise en commun et la recherche de l'intérêt général. Si la réalisation d'une société de partage nécessite un effort de la population pour refuser les biens privés, il est normal que l'acceptation de la monarchie et l'encadrement de celle-ci pour en faire un gouvernement juste passent aussi par des actions des citoyens. Et il est ici question de la dévotion envers le pouvoir. En effet, Ambroise met un point d'orgue non seulement sur l'indulgence des autorités, mais aussi et surtout sur le respect de ces autorités par les citoyens. (se renseigner sur les perses, sans doute un moyen de dire : même les perses et leur monarchie despotique, barbare et violente qui fait des citoyens des pures jouets du roi, ils n'arrivent pas à atteindre le respect que le roi des abeilles obtient par le calme et l'indugence.)L'évêque de Milan met en avant avec les abeilles des principes de sacrifice personnel pour la société, ce qui pousse très loin le principe de respect mutuel et de réciprocité dans l'engagement pour le bien commun. La théorie politique d'Ambroise demande de nouveau un effort général et pas seulement un changement drastique dans la façon d'exercer le pouvoir. En précisant qu'aucun peuple humain ne possède une estime pour son roi aussi grande que les abeilles, il souligne l'un des points majeurs du bon fonctionnement de la monarchie : la confiance ou le rejet qui prend forme derrière les actions des souverains. Dans ce cadre, la nature fournit de nouveau un exemple parfait pour l'évêque, ce qui lui permet de montrer que cet idéal est atteignable si d'autres sociétés y parviennent. Ainsi, si le roi respecte les droits des citoyens, ne cherche pas à s'approprier les biens d'autrui, n'utilise pas la violence sous le coup de la colère, mais œuvre pour la grandeur de la société, alors la population lui doit une obéissance absolue. Ambroise formule également ici une autre de ses pensées politiques importantes : le rejet de la contestation et des luttes civiles pour le pouvoir, point sur lequel nous reviendrons dans le chapitre 2.

\bigskip

Cette dernière citation montre qu'Ambroise peut demander un suivi presque aveugle du souverain, ce qui peut surprendre, mais qui s'appuie sur le fait qu'il est prêt à faire confiance à un gouvernement monarchique, y compris dans sa vision de justice et de clémence, tant que celui-ci respecte l'encadrement qui lui est imposé par la morale et la vertu. Ce fameux contrat moral, qui engage aussi bien les dirigeants que les citoyens, permet à la monarchie de répondre à l'idéal chrétien et de surpasser l'utopie républicaine.

\bigskip

\subsubsection{Questionner les principes de succession et de dynastie}

Pour aboutir à cette dévotion du peuple envers son souverain, Ambroise est conscient de la nécessité d'avoir une confiance absolue envers les institutions et celui qui détient le pouvoir. Si le doute persiste sur les intentions du monarque, il est impossible d'atteindre le niveau de respect mutuel que nous venons de décrire chez les abeilles. Cette exigence fondamntelae oblige l'évêque de Milan à se questionner sur la façon dont il est possible de désigner le dirigeant. Car la présence d'un prince capable d'agir pour le bien commun ne peut pas révéler de la chance, mais bien d'une réflexion autour de la prise du pouvoir et de l'aboutissement au système monarchique. Comme nous l'avons vu, sa théorie du bon gouvernement repose plus sur des fondements moraux et idéaux que sur un changement de structure institutionnel et juridique, il n'exprime donc pas souvent une pensée précise sur le fonctionnement que doit avoir une monarchie concernant les successions. Mais un passage de son \latin{Exameron}, toujours au sujet des abeilles, nous permet entre autre de constater sa critique ferme des principes de dynastie et de tirage au sort qui laissent trop de place au hasard et moisn à la surveillance de la foi et des valeurs. Dans la ruche, l'harmonie qu'il valorise est atteinte sous l'autorité d'un roi qui n'accède pas au pouvoir par les travers humains de l'ambition et de l'héritage :

\begin{quote}
    «~Le roi, en revanche, n’est pas désigné par le sort, car le sort relève du hasard et non du jugement. Il n'est pas non plus désigné par la clameur vulgaire d'une multitude ignorante, qui ne pèse pas les mérites de la vertu ni ne scrute les avantages de l'utilité publique, mais qui vacille dans l'incertitude de son inconstance. Il ne siège pas sur le trône royal en vertu d’un privilège de succession ou de naissance, car celui qui ignore la vie publique ne peut être ni prudent ni sage.\footnote{« \latin{Rex autem non sorte ducitur, quia in sorte eventus est, non iudicium et saepe inrationabili casu sortis melioribus ultimus quisque praefertur, neque inperitae multitudinis uulgari clamore signatur, quae non merita uirtutis expendit nec publicae utilitatis emolumenta rimatur, sed mobilitatis nutat incerto, neque privilegio successionis et generis regalibus thronis insidet, siquidem ignarus publicae conversationis cautus atque eruditus esse non poterit}.~». \cite[\nopp V, 21, 68]{ambroise_csel_32_1}. Traduction personnelle.}~»
\end{quote}

\bigskip

À travers ce passage, Ambroise ne valorise pas seulement la royauté des abeilles, il s'applique à expliquer pourquoi chacun des méthodes d'accession au pouvoir qu'il liste est mauvaise. Son objectif est donc bien de dresser une critique générale des pratiques impériales romaines de son temps, qui reprennent de nombreux éléments qui sont vu dans sa réflexion politique comme des éléments appuyant le risque de l'apparition d'un régime coercitif et tyrannique. Il note en effet qu'à travers la légitimité occtroyée par les outils de la dynastie, de l'acclamation ou du tirage au sort, un problème moral important intervient : comment garantir la justice de la monarchie si le souverain ne possède pas le pouvoir grâce à ses actions et par la reconnaissance des citoyens, mais plutôt par le seul principe d'héritage  dynastique ?



En effet, l'évêque s'efforce de théoriser l'accession au pouvoir en prenant ses distances avec les pratiques impériales romaines de son temps, et plus particulièrement avec le principe de la nomination dynastique de l'héritier (qui se fait dans la très grande majorité des cas en faveur du fils). Il constate que la légitimité automatique conférée par le sang ou le choix arbitraire du prince pose un problème moral et politique majeur : comment garantir qu'un simple héritier sera capable d'incarner et de tenir les vertus exigées par l'exercice du pouvoir ? Par cette remise en question, Ambroise s'inscrit dans le sillage direct de la tradition biblique, où la royauté n'est jamais définitivement acquise par la seule filiation, mais doit être constamment légitimée et réévaluée à l'aune de la foi, de la justice et de la morale. Le monarque chrétien doit faire ses preuves dans tous les aspects définis par la théorie ambrosienne.

\bigskip

Pour illustrer ce rejet des modes de désignation irrationnels ou purement héréditaires, Ambroise recourt une nouvelle fois à l'allégorie de la ruche dans son \latin{Exameron}. Chez les abeilles, l'harmonie de la communauté repose sur un roi qui n'accède pas au pouvoir par les travers humains de l'ambition, du hasard ou de la généalogie :

\bigskip

Ce passage condense une critique systématique des trois grands modes de désignation politique connus du monde gréco-romain. Le meilleur exemple de la vie en communauté, où chacun connaît sa place et fournit son travail, reste celui des abeilles. Parce qu'elles « nomment » leur propre roi selon des critères naturels de mérite, elles peuvent ensuite l'aimer et se donner entièrement à lui pour la réussite de la communauté. À l'inverse, si un dirigeant est choisi par le hasard ou par une simple succession sanguine, le risque de voir émerger des rois incompétents et mauvais s'accroît, entraînant inévitablement des contestations du pouvoir, ce qui représente pour Ambroise le pire des maux pour une société.

\bigskip

La première cible de l'évêque est la désignation « par le sort ». Cette pratique, profondément ancrée dans l'histoire de la République romaine, ne s'est pas totalement éteinte sous l'Empire. Comme le rappelle Julie Bothorel, l'utilisation du tirage au sort (\latin{sortitio}) pour désigner certains postes administratifs, à l'image des proconsuls dans les provinces sénatoriales, se maintient jusqu'à la fin du Haut-Empire\autocite[Voir à ce sujet l'étude de J. Bothorel sur les pratiques de succession provinciale]{bothorel_reference}. Ambroise, ancien gouverneur de province, a une parfaite connaissance de ces rouages institutionnels. En associant le sort à l'absence de jugement et à la ruine de la monarchie, il refuse l'idée qu'un pouvoir absolu puisse reposer sur une dynamique purement aléatoire, incapable de mesurer la vertu de l'individu.

\bigskip

La deuxième critique d'Ambroise vise la « clameur vulgaire d'une multitude ignorante ». Derrière cette formule rhétorique se cache une réalité politique brûlante de l'Antiquité tardive : l'acclamation militaire. L'armée est bien souvent la force décisionnaire qui fait et défait les empereurs par la clameur de ses boucliers. Or, pour Ambroise, cette méthode ne garantit en rien l'« utilité publique » (\latin{publicae utilitatis}) ; elle est surtout le moteur des coups d'État et des usurpations chroniques qui déchirent l'Empire romain\footnote{Nous développerons plus en détail la position d'Ambroise face aux usurpateurs militaires dans la sous-partie [...] du chapitre 2.}. Une monarchie fondée sur la versatilité des légions ne peut aboutir à la stabilité du gouvernement juste qu'il appelle de ses vœux.

\bigskip

Mais c'est la troisième critique qui est la plus audacieuse politiquement : le refus du « privilège de succession ou de naissance ». Si Ambroise ne rejette pas l'institution monarchique en soi, il critique ouvertement ses dérives lorsqu'elle repose sur une logique strictement héréditaire, mère d'orgueil et d'incompétence. Cette réflexion prend tout son sens lorsqu'on la replace dans le contexte constitutionnel romain. En effet, la succession impériale n'a jamais reposé sur des règles de droit parfaitement établies ou sur une loi de primogéniture stricte\autocite[Sur l'absence de constitutionnalité de la succession impériale, voir F. Hurlet]{hurlet_reference}. Le principe dynastique n'est pas un droit fondamental, mais une pratique d'autorité : l'empereur régnant désigne son successeur (le plus souvent en l'associant au pouvoir), une décision qui est ensuite entérinée par l'investiture sénatoriale et l'approbation de l'armée. Le système successoral romain demeure donc un équilibre instable entre ces trois paramètres, rendant la dynastie toujours précaire et soumise à la pression de la contestation.

\bigskip

Cette précarité s'est d'ailleurs transformée en absolutisme à la fin du IIIe siècle. Avec la prise de pouvoir par Dioclétien et la mise en place de la Tétrarchie, l'investiture sénatoriale a perdu sa valeur effective pour aboutir à la seule discrétion du prince régnant : l'empereur choisit un proche (dans sa famille ou son entourage militaire) pour l'élever au rang de César, avant qu'il n'obtienne le statut d'Auguste\autocite[J.-R. Palanque offre une analyse détaillée de cette évolution vers la cooptation absolutiste]{palanque_reference}. C'est précisément ce choix unilatéral que remet en question Ambroise. L'évêque ne condamne pas définitivement et par principe la succession familiale, mais il exige de la soumettre à une stricte conditionnalité morale et religieuse. Pour conforter cette crainte, Ambroise peut d'ailleurs s'appuyer sur la mémoire récente des traumatismes de l'Église, frappée par des successions désastreuses. L'exemple le plus parlant est sans doute celui de l'année 355 : obéissant à une logique strictement dynastique pour consolider son pouvoir, l'empereur Constance II nomme son cousin Julien au titre de César. Une fois parvenu au pouvoir suprême, Julien (dit l'Apostat) se lance dans une vaste entreprise de restauration païenne et de vexation anti-chrétienne. Cet épisode prouve de manière éclatante ce qu'entend Ambroise : l'hérédité, si elle est vidée de la sagesse et de la piété, est une menace mortelle pour l'État et la religion.

\bigskip

On constate toutefois une limite majeure dans la théorisation ambrosienne : bien qu'il déconstruise avec virulence les méthodes de désignation de son temps, il ne propose aucune mesure institutionnelle claire et concrète pour nommer ce « bon roi ». Il reste volontairement évasif, s'en remettant à l'exemple de l'institution par la nature et, de façon plus transcendantale, à l'institution des pouvoirs par Dieu\footnote{L'origine divine du pouvoir temporel et ses implications seront étudiées dans la sous-partie [...].}. Il en ressort que la succession héréditaire ne peut coroborer la réussite morale de la monarchie que si elle est validée par la vertu personnelle du candidat.

\bigskip

Enfin, il est impossible de clore cette réflexion sans souligner le profond paradoxe qui anime l'évêque de Milan. Cette critique acerbe de la dynastie, soupçonnée de pouvoir hisser des incapables ignorants de la chose publique sur le trône d'Auguste, entre en opposition directe avec son pragmatisme et son réalisme politiques. Lorsqu'en 395, l'empereur Théodose Ier meurt, Ambroise n'hésite pas un instant à soutenir et à valoriser la prise de pouvoir de ses deux fils, Honorius et Arcadius. Or, ces derniers, par leur extrême jeunesse et leur inexpérience flagrante, n'ont alors que leur descendance directe avec Théodose comme véritable argument pour régner. Cet épisode prouve qu'Ambroise n'est pas toujours prisonnier de ses propres convictions théoriques. S'il questionne âprement la légitimité des modes de succession pour tenter d'améliorer l'idéal gouvernemental, il possède la finesse d'un homme d'État qui comprend que ses concepts naturalistes restent d'une application périlleuse dans le monde réel, n'hésitant pas à s'accommoder du principe dynastique lorsque la stabilité de l'Empire et la protection de l'Église l'exigent.
